\scanpage{10}
% archive foliation: 9
\textcircled{2.5} Situation où $N$ invariant dans $\mathcal E$ i.e. inv. par $\mathcal G$, et $\mathcal G\to\Gamma$ est surjectif : $\mathcal G'$ est \deleted{\unsure{donc}} muni de suite de composition ayant la structure suivante (\unsure{quots} successifs $\Gamma,\mathbb G,G$)
\[
(24)\quad\diagram{\begin{tikzcd}[ampersand replacement=\&,column sep=large]
G\&\mathbb G\&\Gamma
\arrow[from=1-1,to=1-3,no head,to path={(\tikztostart.north) -- ++(0,.5) -| (\tikztotarget.north)\tikztonodes},"\mathcal G'"]
\arrow[from=1-1,to=1-2,no head,to path={(\tikztostart.south) -- ++(0,-.5) -| (\tikztotarget.south)\tikztonodes},"\mathbb G'"']
\arrow[from=1-2,to=1-3,no head,to path={(\tikztostart.south) -- ++(0,-.5) -| (\tikztotarget.south)\tikztonodes},"\mathcal G"']
\end{tikzcd}}.
\]
L'extension $\mathbb G'$ de $\mathbb G$ par $G=\operatorname{Aut}(\mathcal X'/\mathcal X)$ est de ``nature géométrique'', l'extension de $\Gamma$ par $\mathbb G'$ exprime les \unsure{caractères} arithmétiques de la situation.

Nous nous intéresserons surtout à la situation où $\mathbb G$ et $G$ sont finis, i.e. $\mathbb G'$ fini, et $\Gamma$ infini, i.e. $\bar k$ extension infinie de $k$, \deleted{Considérons} $\mathcal X$ \deleted{\unsure{alors}} (donc aussi $\mathcal X'$) étant de prés. finie sur $k$. [On a encore de plus, $\mathcal X/\bar k/k$ étant fixé, à faire varier $\mathcal X'\inserted{\text{ i.e. }N}$, de façon que $\widetilde{\mathcal X}$ soit la lim. proj. des $\mathcal X'$ i.e. $\bigcap N=1$.] On voit aisément qu'on peut trouver alors un groupe quotient fini $\deleted{\unsure{\Gamma}}\unsure{\gamma_i|\sigma_i}\inserted{(=\Gamma/\Gamma_i)}$ de $\Gamma$
